\chapter{Introduction}
\label{chap:introduction}

Federated Learning (FL) is a distributed training setup where many clients
(e.g.\ phones, hospitals, organizations) collaboratively train one shared
model without pooling their data on a central server --- each client trains
locally on its own private data and only sends back the resulting model
update~\cite{mcmahan2017fedavg}. This is often presented as a
privacy-preserving alternative to centralized training: since raw data
never leaves the client, FL is commonly assumed to keep that data
confidential by construction.\footnote{Code, experiment notebooks, and the
\LaTeX{} source of this report are available at
\url{https://github.com/huynt119/Trustworthy_AI_Project}.}

This project studies whether that assumption actually holds. The gradient
(or weight update) a client sends is not an opaque summary --- it is the
deterministic output of backpropagating a loss on the client's private
batch through a model architecture the server already knows. If that
function can be inverted, the private batch can be reconstructed from the
gradient alone, without any further access to the client. This is the
\emph{gradient inversion} problem, first demonstrated as ``Deep Leakage
from Gradients''~\cite{zhu2019dlg}: given a leaked gradient and a known
model architecture, recover the input/label pair that produced it.

This project focuses on two broad attack strategies: \emph{optimization-based}
attacks, which search for an input/label pair that reproduces the leaked
gradient, and \emph{closed-form} attacks, which solve for it directly from
the per-layer gradient equations, with no search involved. It implements
one representative method from each --- DLG~\cite{zhu2019dlg}/iDLG~\cite{zhao2020idlg}
and Inverting Gradients~\cite{geiping2020inverting} on the optimization
side, R-GAP~\cite{zhu2021rgap} on the closed-form side --- on a common set of
victim architectures, and evaluates both families against a
DP-SGD-style~\cite{abadi2016dpsgd} gradient perturbation defense.

The remainder of this report is organized as follows. Chapter~\ref{chap:related-work}
surveys the work this project builds on. Chapter~\ref{chap:problem-statement}
gives a formal, input/output statement of the FL setting, the threat model,
and the gradient inversion problem. Chapter~\ref{chap:methodology} describes
the attack and defense methods implemented. Chapter~\ref{chap:experiment}
presents the experimental evaluation methodology and results, and
Chapter~\ref{chap:conclusion} concludes with lessons learned.
